\clearpage
\section{مفاهیم پایه}
\subsection{رتینوپاتی نارس و بیماری پلاس}

شبکیه در دوران جنینی شبکه‌ی عروقی خود را کامل می‌کند و تولد زودرس می‌تواند این روند را مختل کند و به رتینوپاتی نارس یا \lr{ROP}\enfoot{بیماری رشد غیرطبیعی عروق شبکیه در نوزاد نارس؛ معادل انگلیسی: \lr{Retinopathy of Prematurity}} بینجامد. داده‌ی این پروژه سه وضعیت \lr{Normal}\enfoot{شبکیه‌ی بدون نشانه‌ی پلاس؛ معادل انگلیسی: \lr{Normal}}، \lr{Pre-Plus}\enfoot{وضعیت میانی میان ظاهر طبیعی و پلاس؛ معادل انگلیسی: \lr{Pre-Plus}} و \lr{Plus}\enfoot{گشادشدگی و پیچ‌خوردگی پیش‌رونده‌ی عروق قطب خلفی؛ معادل انگلیسی: \lr{Plus}} را از هم جدا می‌کند. طبقه‌بندی بین‌المللی ویرایش سوم \lr{ICROP3}\enfoot{مرجع بین‌المللی طبقه‌بندی رتینوپاتی نارس؛ معادل انگلیسی: \lr{ICROP3}} بر پیوستگی این تغییرات و نقش قضاوت متخصص در مرزگذاری تأکید دارد، پس مرز میان \lr{Pre-Plus} و \lr{Plus} پیوسته است و نه دو حالت جدا.

تعریف \lr{Plus} بر شکل رگ استوار است، پس کمیتی که مستقیماً از ماسک عروق اندازه‌گیری می‌شود باید بخشی از آن را در خود داشته باشد. دو راه برای استفاده از این مشاهده وجود دارد؛ اندازه‌گیری صریح از ماسک عروق، یا سپردن همان کار به یک شبکه‌ی تصویری که شکل رگ را از پیکسل می‌آموزد.

\subsection{تصویر فوندوس و نقش آن}

تصویر فوندوس\enfoot{عکس رنگی از سطح داخلی پشت چشم شامل عروق و دیسک بینایی؛ معادل انگلیسی: \lr{fundus image}} مسیر عروق، دیسک بینایی و بافت اطراف را در یک قاب ثبت می‌کند. یک شبکه‌ی تصویری می‌تواند شکل رگ را همراه با رنگ و بافت بیاموزد، بدون آنکه کسی مختصات رگ را به آن بدهد. همین آزادی منبع خطا هم هست، چون تصویر افزون بر بیماری، شرایط ثبت مانند دوربین، نور و کادربندی را نیز حمل می‌کند و مدل می‌تواند بخشی از تصمیم خود را به آن‌ها ببندد.

این خطر در داده‌ی چندمرکزی جدی‌تر می‌شود، چون تصاویر این پروژه از سه منبع مستقل آمده‌اند و شرایط ثبت یکسان نیست؛ اگر یک نشانه‌ی ثبت با توزیع برچسب‌ها همبسته باشد، مدل می‌تواند بدون آموختن بیماری روی داده‌ی آمیخته عدد خوبی بگیرد. به همین دلیل جابه‌جایی بازنمایی‌ها میان منابع بدون استفاده از برچسب هدف اندازه‌گیری شد و سه فولد منبع‌کنارگذاشته ساخته شد که در آن‌ها منبع هدف نه در برازش، نه در پیش‌پردازش و نه در انتخاب مدل دیده نمی‌شود.

\subsection{بیومارکر عروقی}

برای ساختن بیومارکر\enfoot{ویژگی اندازه‌گیری‌شده و قابل‌نام‌گذاری از شکل یا چگالی رگ؛ معادل انگلیسی: \lr{biomarker}}، مدل جداسازی عروق\enfoot{جدا کردن پیکسل‌های رگ از پس‌زمینه‌ی تصویر؛ معادل انگلیسی: \lr{vessel segmentation}} ابتدا احتمال عروقی‌بودن هر پیکسل را محاسبه می‌کند، آستانه‌گذاری این نقشه یک ماسک دودویی می‌سازد و کمیت‌های نهایی از همان ماسک درمی‌آیند. حذف یک رگ باریک یا شمردن یک انعکاس نور به‌عنوان رگ، مقدار بیومارکر را جابه‌جا می‌کند.

این پروژه با پنج نشانگر اسکالر بسته شد؛ چگالی عروق در میدان دید \lr{vessel\_density\_fov}، چگالی خط مرکزی \lr{skel\_density\_fov} و سه بُعد فراکتالی \lr{fractal\_d0}، \lr{fractal\_d1} و \lr{fractal\_d2}. مخرج کمیت‌های چگالی بر پایه‌ی مساحت محتوا تعریف شد، نه مساحت کل کادر؛ شمردن ناحیه‌ی سیاه یا بیرون از میدان دید در مخرج، چگالی را به کادربندی گره می‌زند.

کیفیت این اندازه‌گیری یک محدودیت باز دارد؛ ماسک مرجع متخصص در دسترس نبود، پس صحت جداسازی در برابر مرز انسانی سنجیده نشد و آنچه تأیید شده فقط کامل‌بودن فنی خروجی‌هاست. خود نشانگرها هم وابسته به منبع‌اند و یک نشانگر ثابت در سه مرکز عدد یکسانی نمی‌دهد.

\subsection{شبکه کانولوشنی}

مسیر تصویری این پروژه از یک تصویر رنگی \lr{RGB}\enfoot{سه کانال رنگ سرخ، سبز و آبی؛ معادل انگلیسی: \lr{Red-Green-Blue}} شروع می‌شود و به \lr{EfficientNet-B5}\enfoot{معماری تصویری با مقیاس‌گذاری هم‌زمان عمق، عرض و وضوح؛ معادل انگلیسی: \lr{EfficientNet-B5}} می‌رسد. خروجی این شبکه پیش از طبقه‌بند نهایی یک بردار بازنمایی\enfoot{بردار عددی فشرده‌ای که شبکه از تصویر می‌سازد؛ معادل انگلیسی: \lr{embedding}} است. برخلاف نشانگرها، ابعاد این بردار معنای هندسی نام‌دار ندارند و اطلاعات شکل رگ، رنگ و بافت در همان ابعاد پخش شده است.

همین بردار در دو نقطه به کار می‌رود؛ در یک مدل یک طبقه‌بند خطی روی آن می‌نشیند و در مدل دیگر همان بردار به‌صورت بازنمایی ثابت به \lr{XGBoost}\enfoot{پیاده‌سازی درخت‌های تقویت‌شده با گرادیان؛ معادل انگلیسی: \lr{XGBoost}} داده می‌شود. این تفکیک عملی است؛ با حجم آموزش این مجموعه، تنظیم دقیق کل شبکه بخشی از ظرفیت را صرف برازش نویز می‌کند، در حالی که بازنمایی ثابت تکرارپذیرتر است و مقایسه با نشانگرهای اسکالر را ممکن می‌کند. مسیر عروقی معماری مشابهی دارد؛ نقشه‌ی عروق به \lr{EfficientNet-B4}\enfoot{معماری تصویری با مقیاس‌گذاری هم‌زمان عمق، عرض و وضوح؛ معادل انگلیسی: \lr{EfficientNet-B4}} داده می‌شود و یک بردار بازنمایی عروقی می‌سازد. ورودی این مسیر خودش یک اندازه‌گیری است، نه تصویر خام، پس خطای جداسازی به آن ارث می‌رسد.

\subsection{ترکیب زودهنگام و دیرهنگام}

ساده‌ترین راه استفاده‌ی هم‌زمان از دو نوع اطلاعات، چسباندن آن‌ها پیش از طبقه‌بند است؛ این الگو ترکیب زودهنگام\enfoot{ترکیب ویژگی‌ها پیش از طبقه‌بند نهایی؛ معادل انگلیسی: \lr{early fusion}} نام دارد. الحاق ساده اما یک اشکال ساختاری دارد؛ پنج نشانگر در برابر صدها ستون تصویری سهم کوچکی از ستون‌ها را می‌سازند و یک ستون نشانگری باید در گره‌های درخت از میان ستون‌های تصویری جلو بزند، در حالی که اطلاعات هندسی آن تا حدی در بلوک تصویری هم حاضر است، چون ماسک عروق از همان تصویر ساخته شده است.

ترکیب دیرهنگام\enfoot{ترکیب احتمال یا امتیاز چند مدل مستقل پس از پیش‌بینی؛ معادل انگلیسی: \lr{late fusion}} جای دیگری از همین طیف است؛ امتیاز دو مدل مستقل با یک ضریب وزنی مخلوط می‌شود. در این پروژه ضریب اختلاط روی مرز اشباع شد و نسخه‌ی ترکیب دیرهنگام عملاً به بازنمایی تصویری فروریخت.

ترکیب یادگیرنده پاسخ دیگری به همان اشکال ساختاری بود؛ در مدل‌های ترکیب مشترک، دو رمزگذار تثبیت‌شده هر کدام یک بردار بازنمایی می‌دهند و یک تعامل یادگرفته‌شده آن‌ها را به یک \lr{MLP}\enfoot{شبکه‌ی عصبی چندلایه‌ی متصل؛ معادل انگلیسی: \lr{Multi-Layer Perceptron}} می‌رساند. این طرح در عمل بدتر شد و افزودن نشانگر به همین شبکه هم کاهش پشتیبانی‌شده داد (جدول \autoref{tab:negative}). خطای آموزش این سرها پایین رفت و خطای اعتبارسنجی بالا، و کمینه‌ی خطای اعتبارسنجی در دوره‌ی نخست بود؛ با حجم آموزش این مجموعه، این سرها بیش‌برازش می‌کنند.

قید شرطی کران‌دار آخرین پاسخ آزموده‌شده بود؛ در مدل شرطی‌شده، نشانگرها یک ضریب و یک بایاس کوچک می‌سازند که بازنمایی عروقی را با \lr{FiLM}\enfoot{تعدیل خطی ویژگی با پارامترهای تولیدشده از ورودی شرطی؛ معادل انگلیسی: \lr{Feature-wise Linear Modulation}} بازتنظیم می‌کند. پارامترسازی باقی‌مانده و کران‌دار انتخاب شد تا مقدار اولیه نزدیک همانی بماند و تشخیص‌ها نشان می‌دهند این شرطی‌ساز عملاً همان‌طور ماند؛ جایگزینی هر نشانگر با میانگین آموزشی‌اش تصمیم هیچ نمونه‌ای را عوض نکرد، پس بهبود کالیبراسیون آن نسبت به کنترل هم‌تایش را نمی‌توان به نشانگرها نسبت داد. این نکته به‌عنوان محدودیت طرح تطبیقی در فصل نتایج ثبت شده است.

\subsection{تقسیم گروه‌محور و نشت داده}

یک نوزاد ممکن است در یک معاینه چند تصویر نزدیک به هم داشته باشد؛ اگر یکی در آموزش و دیگری در آزمون بیفتد، مدل ویژگی‌های ثابت همان چشم یا دستگاه را به خاطر می‌سپارد و برآورد عملکرد بالا می‌رود. تقسیم گروه‌محور همه‌ی تصاویر مرتبط را در یک بخش نگه می‌دارد.

جمعیت کانونی این پروژه ۸۸۶۲ مورد کامل در ۴۱۴ گروه است که با تقسیم گروه‌محور به سه بخش آموزش، اعتبارسنجی و آزمون رسید و هیچ گروهی در بیش از یک بخش نیست (جدول \autoref{tab:cohort}). تقسیم یک بار ساخته و با اثر انگشت قفل شد و همه‌ی مدل‌های این گزارش روی همان اجرا شده‌اند. تعریف گروه در سه منبع یکسان نیست؛ در منبع \lr{Plus}\enfoot{بزرگ‌ترین منبع تصویربرداری این مجموعه؛ معادل انگلیسی: \lr{Plus}} و در \lr{FARFUM-RoP}\enfoot{مجموعه‌ی عمومی تصاویر رتینوپاتی نارس؛ معادل انگلیسی: \lr{FARFUM-RoP}} گروه در سطح بیمار و در فارابی در سطح معاینه تعریف شد، چون شناسه‌ی بیمار در آن منبع وجود نداشت. پس در فارابی، نبود هم‌پوشانی معاینه نبود هم‌پوشانی بیمار را اثبات نمی‌کند.

\subsection{معیارهای ارزیابی}

احتمال پیش‌بینی‌شده با یک آستانه به تصمیم تبدیل می‌شود و تغییر آستانه، حساسیت و اختصاصیت را هم‌زمان جابه‌جا می‌کند. منحنی مشخصه‌ی عملکرد گیرنده\enfoot{منحنی بده‌بستان حساسیت و نرخ مثبت کاذب در آستانه‌های مختلف؛ معادل انگلیسی: \lr{ROC}} همین بده‌بستان را در همه‌ی آستانه‌ها نشان می‌دهد و سطح زیر آن\enfoot{سطح زیر منحنی \lr{ROC}؛ خلاصه‌ی توان رتبه‌بندی مدل؛ معادل انگلیسی: \lr{AUC}} یک خلاصه‌ی رتبه‌بندی است، نه دقت تصمیم در یک آستانه. آستانه‌ی تصمیم در این پروژه با شاخص یودن\enfoot{آستانه‌ای که مجموع حساسیت و اختصاصیت را بیشینه می‌کند؛ معادل انگلیسی: \lr{Youden index}} فقط روی اعتبارسنجی انتخاب شد و روی آزمون دوباره تنظیم نشد.

سه معیار دیگر مکمل همین تصویرند. دقت متعادل میانگین بازخوانی سه کلاس است و به نابرابری اندازه‌ی کلاس‌ها حساس نیست، اما مانند \lr{F1} ماکرو\enfoot{میانگین هارمونیک دقت و بازخوانی؛ معادل انگلیسی: \lr{F1 score}} به آستانه وابسته است. امتیاز \lr{Brier}\enfoot{میانگین مربع خطای احتمال پیش‌بینی؛ معادل انگلیسی: \lr{Brier score}} و خطای کالیبراسیون انتظاری \lr{ECE}\enfoot{خطای کالیبراسیون انتظاری؛ معادل انگلیسی: \lr{Expected Calibration Error}} کیفیت خود احتمال را می‌سنجند و کوچک‌تر بودنشان بهتر است. این پنج معیار در هیچ امتیاز ترکیبی ادغام نشده‌اند، چون جهت‌های متفاوت آن‌ها را نمی‌توان ساده جمع کرد.

مجموعه‌ی آزمون این پروژه یک تقسیم استاندارد قفل‌شده است که نتیجه‌ی اصلی روی آن محاسبه شده؛ بنابراین تحلیل‌های بعدی روی همان تقسیم، بازتحلیل اکتشافی ثانویه‌اند و نه یک مجموعه‌ی نگه‌داشته‌شده‌ی تازه برای تأیید. همه‌ی مقایسه‌ها با یک پروتکل واحد و بوت‌استرپ زوجی طبقه‌بندی‌شده سنجیده شده‌اند.
